\bookchapter{附錄F：折疊螢幕選購與使用專題}

\chapterintro{折疊螢幕已經由「未來科技」變成「而家買得到」嘅成熟品類。但佢同直板機喺形態、使用習慣、保養方式上差異好大，值唔值得買、點樣揀、點樣用，呢一篇幫你理清楚。}

\booksection{一、折疊螢幕而家發展到乜嘢程度？}

先講結論：\textbf{折疊螢幕已經唔係「嘗鮮玩具」，但都唔係「人人都要買」嘅主流形態}。佢介乎小眾極客產品同大眾消費品之間，適合特定人羣。

目前主流分兩類：

\begin{itemize}
  \item \textbf{橫向折疊（書本式）}：合上係普通手機大小，展開係接近方形嘅小平板。代表：華為Mate X系列、三星Galaxy Z Fold系列、小米MIX Fold系列、vivo X Fold系列。主打「一部手機=手機+平板」。
  \item \textbf{縱向折疊（翻蓋式）}：合上似粉餅盒，展開係正常尺寸手機。代表：華為Pocket系列、三星Galaxy Z Flip系列、OPPO Find N Flip、小米MIX Flip。主打「小巧便攜+外屏速覽」。
\end{itemize}

2024-2026年嘅關鍵進步：鉸鏈壽命普遍標稱50萬次折疊（每日開合100次可以用13年）、UTG超薄玻璃取代早期塑料膜、摺痕肉眼可見但手感已經好輕微、價格由萬元檔下探到4000-5000元檔。

\begin{tcbinfo}{一句話判斷}
如果你而家用直板機完全夠用、冇「螢幕唔夠大」嘅痛點，就唔好買折疊螢幕。佢係「錦上添花」嘅設備，唔係「雪中送炭」嘅設備。

\end{tcbinfo}

\booksection{二、乜嘢人適合買折疊螢幕？}

根據真實用戶反饋，呢幾類人買折疊螢幕「真香」概率高：

\begin{enumerate}
  \item \textbf{重度閲讀/文檔處理者}：經常喺手機上睇PDF、批註文檔、讀長文。展開後接近7.5英寸嘅內屏，睇A4 PDF唔使縮放，體驗質變。
  \item \textbf{經常出差/開會嘅商務用戶}：唔想帶平板又嫌手機屏細。橫向折疊螢幕展開睇PPT、簽合同、雙屏對照郵件，比直板機效率高一截。
  \item \textbf{鍾意「懸停態」創作嘅人}：折到90度放喺桌上，上半屏取景下半屏撳快門（自拍/延時攝影）、上半屏影片下半屏記筆記、上半屏直播下半屏睇評論——呢個係直板機做到唔到嘅形態。
  \item \textbf{追求便攜又想要大屏嘅女性用戶}：縱向折疊螢幕合上可以塞入細包，展開係正常手機，外屏仲可以當化妝鏡、睇通知。
\end{enumerate}

\begin{tcbwarn}{呢幾類人唔好買}
\begin{itemize}
  \item 手機淨係用嚟刷短影片、傾微信、打電話嘅——折疊螢幕對你毫無增益，反而更重更厚。
  \item 手細或者手腕有腱鞘炎嘅——橫向折疊螢幕展開後普遍230g+，單手握持攰。
  \item 預算緊張、一部手機要用4年以上嘅——折疊螢幕二手保值率低於直板機，維修成本高。
\end{itemize}

\end{tcbwarn}

\booksection{三、選購要點：睇呢6個指標}

唔好俾「折疊」兩個字晃花咗眼，選購嗰陣重點睇：

\begin{itemize}
  \item \textbf{鉸鏈類型}：水滴鉸鏈（摺痕淺、壽命長，目前主流）優於U型鉸鏈（摺痕深、早期方案）。問銷售或查評測確認係邊種。
  \item \textbf{內屏材質}：必須係UTG超薄玻璃，唔好早期嗰種塑料膜。UTG手感接近玻璃，塑料膜指甲劃過會有凹陷感而且易留痕。
  \item \textbf{摺痕表現}：去線下店親手展開睇。2026年嘅合格水平係：亮屏嗰陣正面幾乎睇唔出，斜光下可以睇出一條淺溝，手指劃過有輕微凹陷但唔割手。
  \item \textbf{外屏尺寸}：橫向折疊螢幕嘅外屏越大越好（最好6.4英寸以上），合上可以當正常手機用；縱向折疊螢幕嘅外屏越大越實用（4英寸以上可以回覆消息、睇導航）。
  \item \textbf{重量與厚度}：橫向折疊展開後230-260g係常態，低於230g算優秀；合上厚度12-14mm，低於12mm算薄。縱向折疊合上後可以輕鬆放進恤衫口袋先算合格。
  \item \textbf{售後與保修}：重點睇內屏碎裂嘅維修價格。有啲品牌首年免費換一次內屏，有啲自費要2000-4000元。呢個係折疊螢幕最大嘅隱性成本。
\end{itemize}

\begin{tcbpractice}{線下體驗清單}
\begin{enumerate}
  \item 親手開合10次，感受鉸鏈阻尼係咪順滑、有冇異響。
  \item 展開後斜對住燈光睇摺痕，再正面亮屏睇。
  \item 單手握持展開狀態30秒，感受重量係咪可以接受。
  \item 試用懸停態：折90度放桌上，試下影相同睇影片。
  \item 問清楚內屏單獨維修價格同保修政策。
\end{enumerate}

\end{tcbpractice}

\booksection{四、折疊螢幕獨有嘅使用技巧}

買返嚟之後，呢幾個玩法係直板機體驗唔到嘅：

\booksubsection{分屏與多任務}

橫向折疊螢幕展開後，大部分系統支持上下/左右分屏。典型場景：上半屏睇網課影片，下半屏記筆記；左邊微信右邊文檔；上面地圖導航下面音樂。比直板機分屏實用得多，因為螢幕夠大，每半邊都唔擠。

\booksubsection{懸停態（90度折疊）}

呢個係折疊螢幕嘅「殺手鐧」用法：

\begin{itemize}
  \item \textbf{自拍/合影}：放桌上折90度，上半屏取景，下半屏撳快門，唔使揾人幫手影。
  \item \textbf{延時攝影/夜景}：同樣放桌上穩定拍攝，省三腳架。
  \item \textbf{視像通話}：解放雙手，對方睇你角度都更自然。
  \item \textbf{睇影片}：放桌上當細電視，飛機/高鐵細桌板上特別實用。
  \item \textbf{外屏預覽+內屏拍攝}：用後置攝像頭影自己嗰陣，外屏可以實時顯示畫面，構圖更準。
\end{itemize}

\booksubsection{外屏速覽}

縱向折疊螢幕合上後，外屏可以：睇時間/通知/天氣、回微信（語音轉文字或快捷回覆）、刷公交卡、睇導航、當化妝鏡（調用前置攝像頭）。唔使每次都展開，省電又省事。

\booksubsection{電腦模式}

部分橫向折疊螢幕（華為、三星）展開後可以切換「電腦模式」，外接顯示器鍵鼠後似台簡易PC，臨時處理文檔、回郵件夠用。唔係主力生產力工具，但應急好香。

\booksection{五、保養要點：等折疊螢幕活得耐啲}

折疊螢幕最嬌貴嘅就係內屏同鉸鏈，呢幾條習慣可以顯著延長壽命：

\begin{itemize}
  \item \textbf{唔好喺低温下強行開合}：0℃以下UTG玻璃變脆，開合前先揣喺衫度暖幾分鐘。呢個係官方說明書都會寫但冇人睇嘅重點。
  \item \textbf{避免指甲/硬物接觸內屏}：內屏雖然有保護層但比外屏軟，用指腹操作，唔好用指甲篤。
  \item \textbf{合上嗰陣確認屏內無異物}：一粒沙夾喺中間合上，可能直接壓壞內屏。合上前掃一眼。
  \item \textbf{唔好貼普通鋼化膜}：內屏淨係可以用官方專用膜或UTG保護膜，普通鋼化膜會影響折疊同觸控。建議去官方售後貼膜。
  \item \textbf{鉸鏈唔好入塵液體}：雖然多數標稱IPX8防水，但防塵等級普遍較低。沙灘、工地等環境要小心。
  \item \textbf{開合力度均勻}：唔好單手暴力甩開，雙手由中間勻速展開最穩。
\end{itemize}

\textbf{內屏損壞嘅高發場景}\par

\begin{itemize}
  \item 合上嗰陣夾住耳機線/充電線——直接壓裂內屏。
  \item 放包入面同鑰匙混裝——外屏冇事，內屏俾鑰匙硌出裂痕。
  \item 冬天由凍嘅地方入暖房立刻開合——温差導致玻璃應力變化。
  \item 用咗一年之後覺得「摺痕變深咗」去撕原廠膜——撕膜會連內屏一齊損壞。
\end{itemize}

\booksection{六、折疊螢幕嘅「坑」同心理預期}

買之前心入面要有呢幾本帳：

\begin{itemize}
  \item \textbf{摺痕會一直喺度}：2026年嘅技術做到唔到「完全無痕」，淨係做到「亮屏正面幾乎睇唔出」。如果你係完美主義者，每日盯住摺痕難受，唔好買。
  \item \textbf{更重更厚}：橫向折疊螢幕展開後普遍230g+，單手握持比直板機攰；合上後厚度12-14mm，褲袋入面更鼓。
  \item \textbf{續航普遍一般}：因為要塞鉸鏈機構，電池空間俾壓縮，同價位續航通常唔及直板機。
  \item \textbf{保值率低}：一年後二手價通常係新機嘅50-60\%，比直板機跌得快。
  \item \textbf{維修貴}：內屏總成維修價2000-4000元，鉸鏈維修500-1500元。買之前查清楚保修政策。
  \item \textbf{App適配參差}：主流App已適配大屏，但小眾App可能淨係「拉伸放大」，體驗一般。
\end{itemize}

折疊螢幕唔係「更好嘅手機」，而係「另一種形態嘅手機」。買佢係因為你需要佢獨特嘅形態，唔係因為佢「更先進」。

—— 關於折疊螢幕嘅選購哲學

\booksection{七、主流品牌點樣揀（2026年）}

唔寫具體型號（會過時），淨係寫品牌調性，幫你快速定位：

\begin{itemize}
  \item \textbf{華為Mate X系列}：橫向折疊標杆，鴻蒙生態協同強，外屏比例接近正常手機，合上可以當主力機用。價格高但保值率相對最好。
  \item \textbf{三星Galaxy Z Fold/Flip}：折疊螢幕鼻祖，鉸鏈調校成熟，Z Flip係縱向折疊嘅顏值擔當。國內售後網點較少係短板。
  \item \textbf{小米MIX Fold/Flip}：性價比路線，配置堆料足，價格相對親民。適合想嘗鮮又唔想花太多嘅用戶。
  \item \textbf{vivo X Fold系列}：橫向折疊入面外屏比例做得最接近正常手機嘅，合上當直板機用體驗好。
  \item \textbf{OPPO Find N系列}：橫向折疊入面尺寸偏緊湊嘅，單手可以握持，適合覺得大折疊太重嘅用戶。
  \item \textbf{華為Pocket系列}：縱向折疊入面外屏最大最實用嘅，化妝鏡、防曬檢測等女性向功能做得細緻。
\end{itemize}

選購建議：預算充足追求生態揀華為Mate X；想要最成熟鉸鏈揀三星Z Fold；性價比首選小米/vivo；縱向折疊睇華為Pocket同三星Z Flip。

\begin{tcbinfo}{最後提醒}
折疊螢幕技術每年都喺度迭代，本附錄寫於2026年。買之前務必查最新評測同用戶反饋，重點關注：鉸鏈壽命實測、摺痕表現、內屏維修政策呢三項。線下親手摸永遠比睇評測重要。

\end{tcbinfo}
